% ============================================================================
\clearpage
\section*{v4.8.0-r20 Source-Only: The Tachyon Frame, Complex Rapidity and the Holonomy Test H13}
\addcontentsline{toc}{section}{v4.8.0-r20 Source-Only: The Tachyon Frame, Complex Rapidity and the Holonomy Test H13}
\label{sec:v480-r20-tachyon-complex-rapidity}
% ============================================================================

\paragraph{Purpose.}
Layer r19 found no Lorentz transformation between the sectors of Reading~B and read the earlier reading as a
photon-centred view.  This layer records the pre-registered holonomy test H13 and a post-hoc finding attached to it,
and examines the tachyon frame under the holonomy and spinor structure.  The computations are
\texttt{llm/k35\_test\_h13.py} and \texttt{llm/k36\_tachyon\_frame.py}; only H13 is compared with data.  Earlier layers
are not rewritten.

\noindent\textbf{Source-only status.}
This v4.8.0-r20 layer is source-only.  It does not build a PDF, does not assign a new DOI,
and does not authorize a public release.  The v4.8.0 Version DOI
\href{https://doi.org/10.5281/zenodo.23365244}{10.5281/zenodo.23365244} and the Concept DOI
\href{https://doi.org/10.5281/zenodo.19210728}{10.5281/zenodo.19210728} are unchanged by it.

\subsection*{v4.8.0-r20 the holonomy test H13}
\label{subsec:v480-r20-h13}

The note \texttt{llm/Vorab\_Notiz\_H13\_Wigner\_Holonomie.md} fixed, before any computation, one construction: three
matter frames moving with rapidity \(\eta_H=\operatorname{artanh}(1/\sqrt3)\) in the three triole directions, linked by
pure boosts into a closed loop, with \(\cos\delta\) of the resulting Wigner rotation to be compared with the readout
factor \(f(+)=0.944\pm0.008\) of the backward analysis recorded in \S\ref{subsec:v480-r19-records}.

The holonomy is a pure rotation (boost part \(10^{-16}\)) by \(\delta=0.4985\) rad \(=28.56^\circ\), equal to the
angular defect of the equilateral hyperbolic triangle with vertices at the three velocities.  Then
\(\cos\delta=0.878\), at \(z=-7.9\) from \(f(+)\).  By the fixed criterion H13 is refuted for this construction, and the
holonomy route with triole boosts is closed for \(f(+)\).  Since \(f(+)\) is itself a fit, this removes an explanation
of the factor, not the factor.

\paragraph{Post-hoc finding, not registered and not counting.}
Asked after the result: \(\cos^2(\delta/2)=(1+\cos\delta)/2=0.939\) meets \(f(+)\) at \(z=-0.5\).  It is one of four
obvious functions of \(\delta\) (\(\cos\delta\): \(-7.9\); \(\cos(\delta/2)\): \(+3.1\); \(1-\delta/2\pi\): \(-2.8\)), roughly
the rate expected by chance.  Its parts have different standing.  The half angle is forced once the spinor is read
out: by Theoremlet r06.4 (\S\ref{subsec:v480-r06-null-cone}) \(z\) is quadratic in \(\psi\in\mathbb C^2\), so a rotation
of the frames by \(\delta\) turns \(\psi\) by \(\delta/2\).  The square is a Born-type reading of \(|\langle\psi|U\psi\rangle|^2\).
The orientation is a choice: \(|\langle\psi|U\psi\rangle|^2=\cos^2(\delta/2)+\sin^2(\delta/2)\cos^2\theta\), with \(\theta\) the
angle to the rotation axis, which is normal to the triole plane.  The triole photon \(c(1,\omega,\omega^2)\) is a null
field whose direction \(E\times B\) is \((1,1,1)\), the normal of the plane \(z_1+z_2+z_3=0\); its spinor lies on the
axis and acquires only the phase \(\delta/2=14.28^\circ\).  The factor \(\cos^2(\delta/2)\) belongs to a spinor in the plane,
such as the helicity of a matter frame moving in a triole direction.

\subsection*{v4.8.0-r20 the tachyon frame}
\label{subsec:v480-r20-tachyon}

\paragraph{Theoremlet v4.8.0-r20.1.}
The duality rotation \(z\mapsto iz\) commutes with every element of \(SO(3,\mathbb C)\).  Hence every closed path in
the tachyon sector has the same holonomy as its matter image (for the triole loop \(\delta=28.56^\circ\)), and \(z\) and
\(iz\) have the same two principal null directions.  Matter and tachyon with \(z_T=iz_M\) differ only in the overall
phase of the spinor product.

\paragraph{Proof sketch.}
\(iz\) is a scalar multiple of \(z\), and \(SO(3,\mathbb C)\) acts linearly.  Writing \(z\) as a symmetric 2-spinor
\(F_{AB}\) with \(F_{11}=\tfrac12(z_1-iz_2)\), \(F_{22}=\tfrac12(-z_1-iz_2)\), \(F_{12}=-\tfrac12z_3\) (the polarisation of the
map of r06.4), the principal spinors are the roots of \(F_{11}\lambda^2+2F_{12}\lambda+F_{22}\), which do not change when
\(F\) is multiplied by \(i\).  Checked numerically.  \(\square\)

\paragraph{Theoremlet v4.8.0-r20.2.}
For a speed \(v>1\),
\[
\operatorname{artanh}v=\operatorname{artanh}\frac1v+\frac{i\pi}{2}.
\]
The four-velocity at rapidity \(\eta+i\pi/2\) is \(i\,w\) with \(w=(\sinh\eta,\cosh\eta)\) spacelike: the real
four-velocity turns from timelike to spacelike and the factor \(i\) carries the phase.  For \(v=1/\bsym=1.8203\) the real
part is \(\eta_{\rm sym}=0.6174\) and the imaginary part is \(\pi/2\), the duality angle that relates matter and tachyon
in Theoremlet r07.2 (\S\ref{subsec:v480-r07-duality-angle}).

\paragraph{Proof sketch.}
\(\tanh(x+i\pi/2)=\coth x\); \(\cosh(\eta+i\pi/2)=i\sinh\eta\), \(\sinh(\eta+i\pi/2)=i\cosh\eta\).  \(\square\)

\paragraph{The same structure in three places.}
A phase \(\pi/2\) on the representative and a change of sign of the real invariant occur in the duality rotation
\(z_T=iz_M\), \(s\mapsto-s\) (r07.2), in the superluminal rapidity of r20.2, and in the signed tachyon scale
\(\ell_T^2=-\lambda_T^2\) of \S\ref{sec:r82-tachyon-frame-signed-scale}, where \(\ell_T=i\lambda_T\).  The last is an
allowed effective statement of the release scope; r20 records the parallel and adds nothing to it.

\paragraph{Reading, offered and not adopted.}
The earlier reading places the tachyon at \(-\bsym\), the real part of its rapidity; Reading~B keeps the imaginary part
\(\pi/2\) as the duality angle.  The two readings are then related by keeping or dropping one imaginary part, which is
the photon-centred view of \S\ref{subsec:v480-r19-photon-centred} expressed in rapidities.

\subsection*{v4.8.0-r20 cross terms with complex tachyon rapidity}
\label{subsec:v480-r20-cross-terms}

With the tachyon at \(-\eta+i\pi/2\) the cross terms of \S\ref{subsec:v480-r18-cross-terms} become
\begin{center}
\small
\begin{tabular}{@{}lcc@{}}
\toprule
pair & real (r18) & with \(i\pi/2\)\\
\midrule
matter--photon & \(\cosh\eta=1.197\) & \(1.197\)\\
photon--tachyon & \(\cosh\eta=1.197\) & \(-i\sinh\eta=-0.657\,i\)\\
matter--tachyon & \(\cosh2\eta=1.864=\gL\) & \(-i\sinh2\eta=-1.574\,i\)\\
\bottomrule
\end{tabular}
\end{center}
The matter--tachyon term becomes purely imaginary with modulus \(\sinh2\eta=\sqrt{\gL^2-1}\), the proper velocity of the
pair instead of its Lorentz factor; for tachyons \(\cosh\) and \(\sinh\) exchange their roles.  The real factor \(\gL\)
arises only when the imaginary part is dropped.  This is Theoremlet r19.1 in numbers: Reading~B has no real Lorentz factor
between matter and tachyon.

\paragraph{The tribonacci condition.}
The fixed-point condition \(\mu=1/\bsym(\mu)\) of Theorem~\ref{thm:tribonacci} reads, in rapidities,
\(\cosh2\eta=\coth\eta=\tanh(\eta+i\pi/2)\): the Lorentz factor of the opposite term equals the speed of a tachyon with
real rapidity \(\eta\).  This is a reformulation, not a new value.

\subsection*{v4.8.0-r20 register entries}
\label{subsec:v480-r20-register}

\begin{center}
\small
\begin{tabular}{@{}>{\raggedright\arraybackslash}p{0.34\textwidth}>{\raggedright\arraybackslash}p{0.58\textwidth}@{}}
\toprule
statement & status after r20\\
\midrule
H13, readout factor as triole holonomy & refuted for the fixed construction, \(z=-7.9\) (r20)\\
\(\cos^2(\delta/2)\) for \(f(+)\) & post-hoc, not registered, does not count; half angle from r06.4, square and orientation are readings (r20)\\
tachyon sector of Reading~B & same holonomy and same principal spinors as matter (r20.1)\\
Thm.~\ref{thm:symmetric-frame}, tachyon at \(-\bsym\) & real part of the superluminal rapidity \(\eta_{\rm sym}+i\pi/2\); imaginary part equals the duality angle (r20.2)\\
\(\gL\) as matter--tachyon factor & real only with the imaginary part dropped; with it \(-i\sinh2\eta\) (r20)\\
Thm.~\ref{thm:tribonacci} & equivalent to \(\cosh2\eta=\coth\eta\) (r20)\\
\bottomrule
\end{tabular}
\end{center}

\subsection*{v4.8.0-r20 ledger and guard update}
\label{subsec:v480-r20-ledger-guard}

\[
R_{\rm tach}^{v480\text{-}r20}=(0,0,1)
\]
for (H13 decided, tachyon frame related to both readings through the complex rapidity, physical status of the
imaginary part open).

The layer spends no guard:
\[
R_{\rm glob}^{(v480\text{-}r20,+)}=(0,1,1,1,1),
\qquad
\Gamma_{\rm guard}^{v480\text{-}r20}=(1,1,1,1,1),
\qquad
\beta_4=1/256\ \text{unspent}.
\]
It does not imply exact finite \(C5\) zero closure, exact all-readout terminality,
\(C6/T2\), a global photon-clock/RH theorem, or global all-readout closure.

